    \documentclass{article}
\usepackage{graphicx} % Required for inserting images
\usepackage{amsmath,amsfonts,amssymb, amsthm, mathtools, verbatim, dsfont, hyperref, bbm, tikz, pgfplots}

\usepackage[0T1]{fontenc}
\usepackage{enumitem}
\usepackage{subcaption}
\usepackage{float}
\usepackage{fontspec}
\usepackage{comment}
\setmainfont{TeX Gyre Termes}
\usepackage[
    a4paper,
    left=3cm,
    right=3cm,
    top=3cm,
    bottom=3cm
]{geometry}


\title{Movimento Browniano Fracionário e Aplicações à Matemática Financeira}
\author{Afonso Vilhena Francisco }
\date{Julho de 2026}

\begin{document}

\newtheorem{dfn}{Definição}[section]
\newtheorem{prop}{Proposição}[section]
\newtheorem{thm}{Teorema}[section]
\newtheorem{cor}{Corolário}[section]
\newtheorem{lema}{Lema}[section]

\renewcommand{\proofname}{Demonstração}

\maketitle
\setcounter{section}{-1}

\section{Preliminares}

\subsection{Processos Gaussianos}


\begin{dfn}
Dizemos que a variável aleatória $X$ tem distribuição Gaussiana, ou normal, com valor esperado $\mu \in \mathbb{R}$ e variância $\sigma^2 \geq 0$, e escrevemos $X \sim \mathcal{N}(\mu, \sigma^2)$, se a sua função característica é dada por 
$$\varphi_X(t):=\mathbb{E}[e^{itX}]=e^{it\mu-\frac{1}{2}t^2\sigma^2},\;\;\; \text{ para } t\in \mathbb{R}$$
\end{dfn}

\vspace{0.2cm}

\begin{prop}
As seguintes afirmações são satisfeitas:
\begin{enumerate}[label=\roman*), align=left]

    
    \item Se $\sigma \neq 0$, então $X \sim \mathcal{N}(\mu,\sigma^2)$ tem função densidade $f$ dada por
    $$f_X(x)=\frac{1}{\sqrt{2\pi}\sigma}e^{-{\frac{(x-\mu)^2}{2\sigma^2}}}$$

    \item $X_1 \sim \mathcal{N}(\mu_1,\sigma_1^2)$, $X_2 \sim \mathcal{N}(\mu_2,\sigma_2^2)$ são independentes se e só se $\operatorname{Cov}(X_1,X_2)=0$
    
    \item Se $X_1 \sim \mathcal{N}(\mu_1,\sigma_1^2)$, $X_2 \sim \mathcal{N}(\mu_2,\sigma_2^2)$ são independentes , então
    $$X_1+X_2 \sim \mathcal{N}(\mu_1+\mu_2, \sigma_1^2+\sigma_2^2)$$

    \item Se $a,b\in \mathbb{R}$ e $X \sim \mathcal{N}(\mu, \sigma^2)$, então $aX+b \sim \mathcal{N}(a\mu+b, a^2\sigma^2)$

\end{enumerate}    
\end{prop}

\begin{proof}
\begin{enumerate}[label=\roman*), align=left]
    
    \item 

    
    Vejamos que a função característica é Lebesgue integrável:

    \begin{align*}
        \int_\mathbb{R}|\varphi(t)|dt=\int_\mathbb{R} \left|\exp \left(it\mu-\frac{t^2\sigma^2}{2} \right) \right|dt= \int_\mathbb{R} \exp \left(-\frac{t^2 \sigma^2}{2} \right) dt = \frac{\sqrt{2}}{\sigma}\int_\mathbb{R} \exp(- u^2)du=\frac{\sqrt{2\pi}}{\sigma}<\infty.
    \end{align*}

    \vspace{0.2cm}

    Podemos portanto aplicar o teorema de Levy:

    \begin{align*}
        f_X(x) = \frac{1}{2\pi}\int_\mathbb{R}\varphi(t)e^{-itx}dt=\frac{1}{2\pi}\int_{\mathbb{R}}e^{-it(x-\mu)}e^{-\frac{t^2 \sigma^2}{2}}dt=\frac{\sqrt{2}}{2 \pi \sigma} \int_\mathbb{R}e^{-u^2}e^{-i\frac{\sqrt{2}(x-\mu)}{\sigma}u}du=\frac{1}{\sqrt{2\pi}\sigma}e^{-\frac{(x-\mu)^2}{2\sigma^2}},
    \end{align*}

    onde na última igualdade usamos a seguinte identidade da transformada de fourier de gaussianas
    $$\int_\mathbb{R}e^{-x^2}e^{-iax}dx=\sqrt{\pi}e^{-\frac{a^2}{4}}$$

    \vspace{0.5cm}

    \item A implicação direta é válida para o caso geral, onde $X_1$ e $X_2$ se tratam de variáveis aleatórias de distribuições arbitrárias. De facto, se $X_1 \perp\!\!\!\perp X_2$, então 
    $$\mathbb{E}[X_1X_2]=\mathbb{E}[X_1]\mathbb{E}[X_2],$$ 
    logo
    $$\operatorname{Cov}[X_1,X_2]=\mathbb{E}[X_1X_2]-\mathbb{E}[X_1]\mathbb{E}[X_2]=0.$$
    \vspace{0.1cm}
    Reciprocamente, se $\sigma_{X_1X_2}:=\operatorname{Cov}[X_1,X_2]=0$, então, usando a função característica em dimensão dois, derivada mais abaixo, vem que

    \begin{align*}
        \varphi_{X_1,X_2}(t,s) & =\exp \left( it\mu_{X_1}+is\mu_{X_2}-\frac{1}{2}(t^2\sigma^2_{X_1}+s^2\sigma^2_{X_2}+2ts\sigma_{X_1X_2}) \right) \\
                               & =\exp \left( it\mu_{X_1}+is\mu_{X_2}-\frac{1}{2}(t^2\sigma^2_{X_1}+s^2\sigma^2_{X_2} \right)\\
                               & =\exp \left( it\mu_{X_1} -\frac{1}{2}t^2\sigma^2_{X_1}\right) \exp \left( is\mu_{X_2} -\frac{1}{2}s^2\sigma^2_{X_2} \right)
                               =\varphi_{X_1}(t)\varphi_{X_2}(s)
    \end{align*}

    \vspace{0.5cm}

    \item Usando a definição de função característica,
    \begin{align*}
        \varphi_{X_1+X_2}(t)&=\mathbb{E}[e^{it(X_1+X_2)}]=\mathbb{E}[e^{itX_1}]\mathbb{E}[e^{itX_2}]=\varphi_{X_1}(t)\varphi_{X_2}(t)\\
                            &=\exp\left(it\mu_1-\frac{t^2\sigma_1^2}{2} \right) \exp \left( it\mu_2-\frac{t^2\sigma_2^2}{2} \right)\\
                            &=\exp\left( it(\mu_1+\mu_2) -\frac{1}{2}t^2\left( \sigma_1^2+\sigma_2^2 \right) \right).
    \end{align*}

    Portanto, $$X_1+X_2 \sim \mathcal{N}(\mu_1+\mu_2, \sigma_1^2+\sigma_2^2)$$

    \vspace{0.5cm}

    \item Mais uma vez, pela definição de função característica,
    \begin{align*}
        \varphi_{aX+b}(t)& = \mathbb{E}[e^{it(aX+b)}]=\mathbb{E}[e^{i(at)X}]\mathbb{E}[e^{itb}]=\varphi_X(at)e^{itb}\\
                         & = \exp \left( i(at)\mu-\frac{1}{2}(at)^2\sigma^2 \right) \exp(itb)\\
                         & = \exp \left( it(a\mu+b)-\frac{1}{2}t^2(a\sigma)^2 \right).
    \end{align*}

    Portanto, $$aX+b \sim \mathcal{N}(a\mu+b, a^2\sigma^2)$$

    
\end{enumerate}

    
\end{proof}

\vspace{1cm}

\begin{dfn}[Vetor Gaussiano]
Um vetor aleatório $\textbf{X}=(X_1,...,X_n)$ diz-se ter distribuição gaussiana n-dimensional se toda a combinação linear das suas componentes tem distribuição gaussiana. Isto é, se
$$\text{Para qualquer } \mathbf{t}\in \mathbb{R}^n, \text{a variável aletória } Y=\mathbf{t}^T \textbf{X} \text{ tem distribuição normal}$$
Nesse caso, dizemos que $\mathbf{X}$ é um vetor aleatório Gaussiano, ou alternativamente, que $X_1,...,X_n$ tem distribuição Gaussiana conjunta.    
\end{dfn}

\vspace{1cm}

É imediato ver que se  $\mathbf{\mu}=(\mu_1,...,\mu_n)$, onde $\mu_i=\mathbb{E}[X_i]$, para $i=1,...,n$ e $\Sigma=[\operatorname{Cov(X_i,X_j)}]_{i,j \in {1, ..., n}}$, então $$Y \sim \mathcal{N}(\mathbf{t}^\top\boldsymbol{\mu}, \mathbf{t}^{\top}\Sigma \mathbf{t}).$$

\vspace{0.5cm}
De facto, $\mathbb{E}[Y]=\mathbb{E}[\mathbf{t}^\top \mathbf{X}]=\mathbf{t}^\top \mathbb{E}[\mathbf{X}]=\mathbf{t}^\top \boldsymbol{\mu}$ e

$\operatorname{Var[Y]}=\operatorname{Var[\mathbf{t}^\top \mathbf{X}]}=\sum_{i,j=1}^n{t_it_j \operatorname{Cov(X_i,X_j)}}=\mathbf{t}^\top \Sigma \mathbf{t}$

\vspace{0.5cm}
Como consequência da observação anterior, os vetores Gaussianos são completamente caracterizados pelo seu valor médio e pela sua matriz de covariância. De facto, a sua função característica é dada por 

$$\varphi_X(t)=\mathbb{E}[e^{it^T\mathbf{X}}]=\mathbb{E}[e^{iY}]=\varphi_Y(1)=e^{i\mathbf{t}^\top\boldsymbol{\mu}-\frac{1}{2}\mathbf{t}^\top \mathbf{\Sigma} \mathbf{t}}$$

\vspace{1cm}

\begin{dfn}
Dizemos que um processo estocástico $\{X_t;\ t \in T\}$ é gaussiano se para qualquer número finito de índices $t_1,...,t_n \in T$ o vetor aleatório $\mathbf{X}=(X_{t_1},..., X_{t_n})$ é Gaussiano. 
No caso em que $\mu(t)=\mathbb{E}[X_t]\equiv0$, dizemos que $\{X_t;\ t \in T\}$ é um processo gaussiano centrado.
\end{dfn}

\vspace{2cm}

\subsection{Movimento Browniano}

Considere-se a experiência de lançar uma moeda justa, isto é, a probabilidade de sair qualquer uma das faces é $\frac{1}{2}$, repetidamente. Denotando por $w_j$ o resultado o j-ésimo lançamento, definimos

\begin{equation*}
    X_j= \begin{cases}
        \;\;\,1, \,\,\,\;\;\;w_j=Cara\\
        -1, \;\;\;\,\; w_j=Coroa
    \end{cases}
\end{equation*}

\noindent Além disso, definimos o contador até ao lançamento $n$ como

\begin{equation*}
    M_n= \sum_{i=0}^nX_i,
\end{equation*}
com $M_0=0$. O processo $(M_n)_{n \in \mathbb{N}_0}$ é chamado \textit{passeio aleatório simétrico}.

\vspace{0.5cm}

O passeio aleatório simétrico está definido apenas para  observações discretas. Estamos interessados em obter o contínuo análogo do passeio aleatório. Para isso \textit{acelera-se} a escala do tempo e \textit{diminui-se} o tamanho do passo do passeio aleatório. Rigorosamente, fixado um $n \in \mathbb{N}$ definimos o \textit{passeio aleatório simétrico normalizado }
\begin{equation*}
    W^{(n)}(t) = \frac{1}{\sqrt{n}}M_{nt}
\end{equation*}

sempre que $nt$ for um inteiro.

\vspace{1cm}

No limite $n \rightarrow \infty$ obtemos o movimento Browniano. Muitas das propriedades do movimento Browniano são herdadas do passeio aleatório normalizado, o que leva à seguinte definição.

\vspace{0.5cm}

\begin{dfn}
    Um movimento Browniano é um processo estocástico $(W_t)_{t \geq 0}$ com as seguintes propriedades:
    \begin{enumerate}
        \item $W_0=0$;
        \item A função $t \mapsto W_t$ é contínua, quase certamente;
        \item O processo $(W_t)_{t \geq 0}$ tem incrementos estacionários e independentes;
        \item O incremento $W_{t_{i+1}}-W_{t_i} \sim \mathcal{N}(0,t_{i+1}-t_i)$;
    \end{enumerate}
\end{dfn}

\vspace{1cm}

Como os incrementos 
\begin{equation*}
    W_{t_1},W_{t_2}-W_{t_1},...,W_{t_n}-W_{t_{n-1}}
\end{equation*}
são independentes e normalmente distribuidos, $(W_t)_{t \geq 0}$ é um processo Gaussiano, e portanto completamente caracterizado pela sua função média e função covariância.\\

Seja $t\geq 0$

\begin{align*}
    \mathbb{E}[W_t]=\mathbb{E}[W_t-W_0]=0
\end{align*}

Dados $0 \leq s <t$ 
\begin{align*}
    \mathbb{E}[W_sW_t]&=\mathbb{E}[W_s(W_t-W_s)+W_s^2]\\
                      &=\mathbb{E}[W_s]\mathbb{E}[W_t-W_s]+\mathbb{E}[W_s^2]\\
                      &=0 + \operatorname{Var}[W_s]\\
                      &=s\\
                      &=\min(s,t)
\end{align*}


\vspace{0.5cm}

Mais especificamente, o movimento Browniano é um processo Gaussiano centrado.

\vspace{0.5cm}

\begin{prop}
    O movimento Browniano é uma martingale~\ref{dfn:martingale}.
\end{prop}

\begin{proof}
    Sejam $0 \leq s <t$

    \begin{align*}
        \mathbb{E}[W_t |\mathcal{F}_s]&=\mathbb{E}[(W_t-W_s)+W_s|\mathcal{F}_s]\\
                                      &=\mathbb{E}[W_t-W_s|\mathcal{F}_s]+\mathbb{E}[W_s|\mathcal{F}_s]\\
                                      &=\mathbb{E}[W_t-W_s]+W_s\\
                                      &=0+W_s=W_s\;,
    \end{align*}
    onde na  terceira igualdade usámos o facto de $W_t-W_s$ ser independente de $W_s$ e o facto de $W_s$ ser $\mathcal{F}_s$-mensurável.
\end{proof}

\vspace{0.5 cm}

\begin{dfn}
    Seja $(X_t)_{t \geq 0}$ um processo estocástico. Definimos variância quadrática de $(X_t)_{t \geq 0}$ como
    $$[X]_t=\lim_{\|P\| \rightarrow 0}\sum_{k=1}^n(X_{t_k}-X_{t_{k-1}})^2,$$

    onde $P$ percorre todas as partições do intervalo $[0,t]$.
\end{dfn}

\vspace{0.5cm}

\begin{prop}
    O movimento Browniano, $(W_t)_{t \geq 0}$ tem variância quadrática igual a $t$.
\end{prop}


\begin{proof}
    Considere-se uma partição  $\Pi=\{t_0,t_1,...,t_n\}$ de $[0,t]$, isto é, $t_0=0,t_n=t$.
    
    Defina-se $$Q_\Pi= \sum_{i=0}^{n-1}(W_{t_{i+1}}-W_{t_i})^2.$$

    Queremos ver que $\lim_{||\Pi|| \rightarrow 0}Q_\Pi=t$.

    Comecemos por notar que

    \begin{equation*}
        \mathbb{E}[(W_{t_{i+1}}-W_{t_i})^2]=\operatorname{Var}[W_{t_{i+1}}-W_{t_i}]=t_{i+1}-t_i,
    \end{equation*}
    logo 
    \begin{equation*}
        \mathbb{E}[Q_\Pi]= \sum_{i=0}^{n-1}\mathbb{E}[(W_{t_{i+1}}-W_{t_i})^2]=\sum_{i=0}^{n-1}t_{i+1}-t_i=t.
    \end{equation*}

    Além disso,

    \begin{align*}
        \operatorname{Var}[(W_{t_{i+1}}-W_{t_i})^2]&=\mathbb{E}[(W_{t_{i+1}}-W_{t_i})^4]-\mathbb{E}^2[(W_{t_{i+1}}-W_{t_i})^2]\\
                                                   &=\mathbb{E}[(W_{t_{i+1}}-W_{t_i})^4]-(t_{i+1}-t_i)^2\\
                                                   &=3(t_{i+1}-t_i)^2-(t_{i+1}-t_i)^2\\
                                                   &=2(t_{i+1}-t_i)^2,
    \end{align*}
    logo
    \begin{align*}
        \operatorname{Var}[Q_\Pi]&=\sum_{i=0}^{n-1}\operatorname{Var}[(W_{t_{i+1}}-W_{t_i})^2]\\
                                 &=\sum_{i=0}^{n-1}2(t_{i+1}-t_i)^2\\
                                 &\leq \sum_{i=0}^{n-1}2||\Pi||(t_{i+1}-t_i)=2||\Pi||t                                 
    \end{align*}
    o que implica que $\lim_{||\Pi|| \rightarrow 0}\operatorname{Var}[Q_\Pi]=0$.
    \vspace{0.5cm}

    Note-se que $$||X_n-X||_2=\left(\int|X_n-X|^2d\mathbb{P} \right)^\frac{1}{2}=\mathbb{E}[(X_n-X)^2]^\frac{1}{2}$$
    
    Portanto, $$X_n \overset{L^2}{\rightarrow}X \Leftrightarrow \mathbb{E}[(X_n-X)^2] \overset{n \rightarrow \infty}{\longrightarrow}0$$
    
    Mas,
    \begin{align*}
        \operatorname{Var}(Q_\Pi)&=\mathbb{E}[(Q_\Pi-\mathbb{E}[Q_\Pi])^2]\\
                                 &=\mathbb{E}[(Q_\Pi-t)^2] \overset{||\Pi|| \rightarrow 0}{\longrightarrow} 0
    \end{align*}

    Portanto, $$Q_\Pi \overset{L^2}{\longrightarrow}t$$

    E como convergência em $L^2$ implica convergência em medida, temos que 

    $$Q_\Pi \overset{\mathbb{P}}{\longrightarrow}t$$
\end{proof}

\begin{prop}
    O movimento Browniano não é diferenciável em nenhum ponto, quase certamente.
\end{prop}
\begin{proof}
    A prova pode ser consultada em \cite{Rostek2009}
\end{proof}


\vspace{1cm}
\subsection{Cálculo de Itô}

O objetivo desta secção é definirmos o integral de Itô: fixado t queremos definir
\begin{equation*}
    \int_0^t\Delta(s)dW(s),
\end{equation*}
onde $(W_t)_{t\geq 0}$ é um movimento Browniano com uma filtração $(\mathcal{F}_t)_{t \geq 0}$ e a integranda $\Delta(t)$ um processo estocástico adaptado.

Poderia-se tentar definir o integral da seguinte forma
\begin{equation*}
    \int_0^t\Delta(s)dW(s)=\int_0^t\Delta(s)W'(s)ds,
\end{equation*}

onde o termo da direita é apenas um integral de Lebesgue com respeito ao tempo. Acontece que, como visto anteriormente, o movimento Browniano não é diferenciável, logo o termo $W'(s)$ nem sequer está definido. Portanto, não podemos proceder deste modo.

\vspace{0.5cm}

Ao invés, iremos primeiro definir o integral de Itô para integrandas simples e depois estender para integrandas não simples como o limite do integral de integrandas simples. 

Mais precisamente, seja $\Pi=\{t_0,t_1,...,t_n\}$ uma partição de $[0,t]$ e $(X(s))_{s \geq 0}$ um processo adaptado à filtração gerada pelo movimento Browniano, $(\mathcal{F}_s^W)_{s \geq 0}$.

Adicionalmente, assumimos que,
$$\mathbb{E} \left[ \int_0^tX^2(s)ds \right]< \infty$$
e que se trata de um processo simples, isto é,
$$X(s)=\sum_{i=0}^{n-1}X_i \mathds{1}_{[t_i,t_{i+1}]}$$


onde $X_i$ são variáveis aleatórias limitadas e  $\mathcal{F}^W_{t_i}$-mensurável.

Definimos então o integral de Itô para $(X(s))_{s \geq 0}$ como

\begin{equation*}
    \int_0^tX(s)dW(s) \coloneq \sum_{i=0}^{n-1} X_i[W_{t_{i+1}}-W_{t_i}]
\end{equation*}

\vspace{0.5cm}

Seja agora $(X(s))_{s\geq 0}$ um processo geral, não necessariamente simples, adaptado à filtração gerada pelo movimento Browniano. Então existe uma sequência de processos simples, $\{X_n(s), s \geq 0\}_{n \in \mathbb{N}}$ tal que $||X_n(s)-X(s)||_2 \rightarrow 0$. Podemos então definir o integral de Itô como
$$\int_0^tX(s)dW(s):=\lim_{n \rightarrow\infty} \int_0^tX_n(s)dW(s),\;\;\;\; \text{para } 0\leq t \leq T.$$
    

\begin{thm}
    Seja $f(t,x)$ uma função para a qual as derivadas parciais $\frac{\partial f}{\partial x},\frac{\partial f}{\partial t},\frac{\partial^2 f}{\partial x^2}$ estão definidas e são continuas. Então, para todo o $T \geq 0$,
    \begin{align*}
        f(T,W(T))=f(0,W(0))+\int_0^T\frac{\partial f}{\partial t}(t,W(t))dt+\int_0^T\frac{\partial f}{\partial x}(t,W(t))dW(t)+\frac{1}{2}\int_0^T\frac{\partial^2 f}{\partial x^2}(t,W(t))dt
    \end{align*}
\end{thm}
\begin{proof}
    Um prova deste facto pode ser encontrada em \cite{Shreve2004}.
\end{proof}

\vspace{0.2cm}

Podemos escrever a fórmula de Itô na forma diferencial,

$$\boxed{df(t,W(t))=\frac{\partial f}{\partial t}(t,W(t))dt+\frac{\partial f}{\partial x}(t,W(t))dW(t)+\frac{1}{2}\frac{\partial^2 f}{\partial x^2}(t,W(t))dt}$$

\newpage


\section{Movimento Browniano Fracionário e propriedades básicas}

Introduzido em 1940 por \textit{Andrey Kolmogorov} e mais tarde popularizado por \textit{Mandelbrot \& Van Ness}, o movimento Browniano fracionário (\textit{mBf}) é uma generalização do movimento Browniano. \\



\noindent Nesta secção definimos o \textit{mBf} e estudamos algumas das suas propriedades mais importantes.

\vspace{0.5cm}

\begin{dfn}
Seja $H \in(0,1)$. Um movimento Browniano fracionário (mBf) de parâmetro de Hurst $H$ é um processo Gaussiano centrado $\{B_t^H, t\geq0\}$ com função de covariância 

\begin{equation*}
    \mathbb{E}[B_t^HB_s^H]=\frac{1}{2}(t^{2H}+s^{2H}-|t-s|^{2H})
\end{equation*}

\end{dfn}

Como visto anteriormente, um processo Gaussiano depende apenas da sua função média e de covariância. Ora sendo mBf um processo centrado, temos que 

\begin{equation*}
    \mathbb{E}[B_t^H]=0, \text{ para qualquer } t>0
\end{equation*}
\begin{equation*}
    \operatorname{Cov}(B_s^H,B_t^H)=\mathbb{E}[B_s^HB_t^H]-\mathbb{E}[B_s^H]\mathbb{E}[B_s^H]=\mathbb{E}[B_s^HB_t^H],
\end{equation*}

portanto a definição acima caracteriza completamente o processo estocástico.

\vspace{0.5cm}

Vejamos que o processo definido acima existe de facto. Para isso introduzimos a noção de função \textit{positiva definida}.

\vspace{0.5cm}

\begin{dfn}
    Uma função $f:Y^2 \rightarrow W$ diz-se \textbf{simétrica} se para quaisquer $y_1,y_2 \in Y$, $$f(y_1,y_2)=f(y_2,y_1).$$
    
    Uma função $f: T^2 \rightarrow \mathbb{R} $ simétrica diz-se \textbf{definida positiva} se para qualquer $n \in \mathbb{N}$ e  quaisquer $x_1,..,x_n \in \mathbb{R}$
    \begin{equation*}
        \sum_{i,j}^nx_ix_jf(t_i,t_j) \geq 0
    \end{equation*}
\end{dfn}

\vspace{0.5cm}
Podemos então enunciar o seguinte resultado
\vspace{0.5cm}

\begin{thm}[Kolmogorov] Seja $\Gamma:T^2 \rightarrow \mathbb{R}$ é uma função simétrica. Então existe um processo Gaussiano centrado com função covariância $\Gamma$ se e só se $\Gamma$ for positiva definida.
\end{thm}
\begin{proof}
    A prova pode ser consultada em \cite{RevuzYor1999}.
\end{proof}
\vspace{0.5cm}

\begin{prop}
A função $\Gamma(s,t)=\frac{1}{2}(t^{2H}+s^{2H}-|t-s|^{2H})$ é positiva definida para $H\in (0,1]$
\end{prop}
\begin{proof}
    A prova pode ser encontrada em \cite{Nourdin2012}
\end{proof}
\vspace{0.5cm}

\noindent De uma aplicação direta do Teorema 1.1 segue o resultado pretendido.

\vspace{0.5cm}

\begin{cor}O processo definido anteriormente como movimento Browniano fracionário existe.
\end{cor}

\vspace{1cm}

O movimento Browniano fracionário é uma generalização do movimento Browniano no sentido em que para $H=\frac{1}{2}$ recuperamos o movimento Browniano. Como referido em \cite{Sznitman2013}, um processo estocástico é um movimento Browniano se e só se for um processo gaussiano centrado de função de covariância $\Gamma(s,t)=\min(s,t)$. De facto, como o movimento browniano fracionário é um processo gaussiano centrado, o processo resultante de $H=\frac{1}{2}$ também o é, e temos
 
$$\mathbb{E}[B_t^{1/2}B_s^{1/2}]=\frac{1}{2}(t+s-|t-s|)=\min(t,s).$$



\vspace{1cm}

\subsection{Incrementos}

\begin{thm}
Para $H \in (0,1)$  o mBf tem incrementos estacionários, isto é, para quaisquer $t,s\geq 0$
$$B^H_{t+s}-B^H_t \overset{\mathcal{D}}{=}B^H_s$$
\end{thm}
\begin{proof}
Ora, $$\mathbb{E}[B^H_{t+s}-B^H_t]=\mathbb{E}[B^H_{t+s}]-\mathbb{E}[B^H_t]=0-0=0$$
e
\begin{align*}
    \operatorname{Var}[B^H_{t+s}-B^H_t]&=\operatorname{Var}[B^H_{t+s}]+\operatorname{Var}[B^H_{t}]-2\operatorname{Cov}[B^H_{t+s},B^H_t]\\
    &=\mathbb{E}[(B^H_{t+s})^2]+\mathbb{E}[(B^H_t)^2]-2\mathbb{E}[B^H_{t+s}B^H_t]\\
    &=\frac{1}{2}((t+s)^{2H}+(t+s)^{2H}-0)+\frac{1}{2}(t^{2H}+t^{2H}-0)-2\times\frac{1}{2}((t+s)^{2H}+t^{2H}-s^{2H})\\
    &=s^{2H}=\operatorname{Var}[B_s^H]
\end{align*}
\end{proof}

\begin{comment}
\begin{proof}
    \begin{align*}
        \mathbb{E}[(B^H_{t_1-s_1}-B^H_{t_1})(B^H_{t_2-s_2}-B^H_{t_2})] &= \mathbb{E}[B^H_{t_1-s_1}B^H_{t_2-s_2}]-\mathbb{E}[B^H_{t_1-s_1}B^H_{t_2}]-\mathbb{E}[B^H_{t_1}B^H_{t_2-s_2}]+\mathbb{E}[B^H_{t_1}B^H_{t_2}]\\
        &=\frac{1}{2}((t_1-s_1)^{2H}+(t_2-s_2)^{2H}-|t_1-s_1-t_2+s_2|^{2H}\\
        &-(t_1-s_1)^{2H}-t_2^{2H}+|t_1-s_1-t_2|^{2H}\\
        &-t_1^{2H}-(t_2-s_2)^{2H}+|t_1-t_2-s_2|^{2H}\\
        &+t_1^{2H}+t_2^{2H}-|t_1-t_2|^{2H})\\
        &=\frac{1}{2}(|t_1-s_1-t_2|^{2H}+|t_1-t_2-s_2|^{2H}-|t_1-s_1-t_2+s_2|^{2H}-|t_1-t_2|^{2H})
    \end{align*}
\end{proof}
\end{comment}

\vspace{0.5cm}

Verifiquemos agora a dependência dos incrementos do mBf. Sejam $0<s_1<t_1<s_2<t_2$, temos que 

\begin{align*}
    \mathbb{E}[(B^H_{t_1}-B^H_{s_1})(B^H_{t_2}-B^H_{s_2})]&=\mathbb{E}[B^H_{t_1}B^H_{t_2}]-\mathbb{E}[B^H_{t_1}B^H_{s_2}]-\mathbb{E}[B^H_{s_1}B^H_{t_2}]+\mathbb{E}[B^H_{s_1}B^H_{s_2}]\\
    &=\frac{1}{2}(t_1^{2H}+t_2^{2H}-|t_1-t_2|^{2H}-t_1^{2H}-s_2^{2H}+|t_1-s_2|^{2H}-s_1^{2H}-t_2^{2H}+|s_1-t_2|^{2H}\\
    &\hspace{0.4cm}  +s_1^{2H}+s_2^{2H}-|s_1-s_2|^{2H})\\
    &=\frac{1}{2}(|t_1-s_2|^{2H}-|t_1-t_2|^{2H}+|s_1-t_2|^{2H}-|s_1-s_2|^{2H})\\
    &=\frac{1}{2}((s_2-t_1)^{2H}-(t_2-t_1)^{2H}+(t_2-s_1)^{2H}-(s_2-s_1)^{2H})
    \begin{cases}
        >0, \text{ se } H >1/2 \\
        =0, \text{ se } H =1/2 \\
        <0, \text{ se } H<1/2
    \end{cases}
\end{align*}
onde a última afirmação segue da análise do comportamento de $f(x)=x^{2H}$ em função de $H$.

\vspace{0.5cm}

Intuitivamente, para $H>1/2$, se o movimento \textit{sobe} (resp. \textit{desce}) num certo intervalo há maior probabilidade de \textit{subir} (resp. \textit{descer}) num intervalo subsequente disjunto. Nesse caso, dizemos que o \textit{mBf} possui incrementos \textit{persistentes}.

Análogamente, para $H<1/2$, se o movimento \textit{sobe} (resp. \textit{desce}) num certo intervalo há maior probabilidade de \textit{descer} (resp. \textit{subir}) num intervalo subsequente disjunto. Nesse caso, dizemos que o \textit{mBf} possui incrementos \textit{antipersistentes}.

Finalmente, para $H=1/2$, temos incrementos independentes, isto é, o movimento em intervalos anteriores não afeta o movimento em intervalos subsequentes, desde que disjuntos. Esta é uma propriedade característica do movimento Browniano.


\begin{figure}[H]
    \centering

    \begin{subfigure}{0.48\textwidth}
        \centering
        \includegraphics[width=\textwidth]{persistentes.pdf}
        \caption{$H=0.8$}
    \end{subfigure}
    \hfill
    \begin{subfigure} {0.48\textwidth}
        \centering
        \includegraphics[width=\textwidth]{antiperssistentes.pdf}
        \caption{$H=0.3$}
    \end{subfigure}


    \begin{subfigure} {0.75\textwidth}
        \centering
        \includegraphics[width=\textwidth]{independentes.pdf}
        \caption{$H=0.5$}
    \end{subfigure}

    \caption{Incrementos do \textit{mBf}}
\end{figure}


\vspace{0.5cm}


\subsection{Autossimilaridade e  inversão de tempo}

A propriedade de autossimilaridade afirma intuitivamente que se \textit{comprimirmos} o eixo do tempo por um fator de $a$, e \textit{dilatarmos} o eixo vertical por um fator de $a^H$ o \textit{mBf} resultante terá uma aparência "similar" ao movimento original. Dizemos que o movimento Browniano fracionário admite uma propriedade de \textit{autossimilaridade}. Rigorosamente,

\vspace{0.5cm}

\begin{thm}
Seja $a>0$, então $(B_{at}^H)_{t\geq0}\overset{\mathcal{D}}{=}(a^HB_t^H)_{t\geq0}$ 
\end{thm}

\begin{proof}
Basta ver que as funções de covariância de $(B_{at}^H)_{t\geq0}$ e de $(a^HB_t^H)_{t\geq0}$ coincidem. Ora,
\begin{align*}
    \mathbb{E}[(a^HB_t^H)(a^HB_s^H)]&=a^{2H}\mathbb{E}[B_t^HB_s^H]\\
    &=\frac{a^{2H}}{2}(t^{2H}+s^{2H}-|t-s|^{2H})\\
    &=\frac{1}{2}[(at)^{2H}+(as)^{2H}-|at-as|^{2H}]=\mathbb{E}[B_{at}^HB_{as}^H]
\end{align*}
\end{proof}

\begin{figure}[H]
    \centering
    \includegraphics[width=1\linewidth]{0fe13f67-11c4-4eec-a70e-b10350bc1a42.pdf}
    \caption{Autossimilaridade do movimento Browniano fracionário (H=0.7)}
    \label{fig:placeholder}
\end{figure}

\vspace{1cm}

\noindent Outra propriedade de similaridade do \textit{mBf} é a de inversão de tempo.

\begin{thm}
    Para $t\neq 0$ temos,
    \begin{equation*}
        t^{2H}B^H_{1/t} \overset{\mathcal{D}}{=}B^H_t
    \end{equation*}
\end{thm}
\begin{proof}
    \begin{align*}
       \mathbb{E}[t^{2H}B^H_{1/t}s^{2H}B^H_{1/s}]&=(ts)^{2H}\mathbb{E}[B^H_{1/t}B^H_{1/s}] \\
       &=(ts)^{2H}\frac{1}{2}[(1/t)^{2H}+(1/s)^{2H}-|1/t-1/s|^{2H}]\\
       &=\frac{1}{2}[s^{2H}+t^{2H}-|s-t|^{2H}]=\mathbb{E}[B^H_tB^H_s]
    \end{align*}
\end{proof}


\subsection{Continuidade}
Queremos analisar a continuidade do movimento Browniano fracionário. Isto é, fixada uma realização $\omega$, analisar a continuidade da função
$t \mapsto B_t^H(\omega)$. Para tal, introduzimos a definição de modificação de um processo estocástico e de seguida invocamos um resultado famoso: o \textit{teorema da continuidade de Kolmogorov-Chentsov}.

\vspace{1cm}

\begin{dfn}
    Sejam $\{X_t, t \in T\},\{Y_t, t \in T\}$ dois processos estocásticos definidos no mesmo espaço de probabilidade $(\Omega, \mathcal{F}, \mathbb{P})$, com  os mesmos conjuntos de índices $T$ e conjunto amostral $S$. O processo $Y$ diz-se uma modificação de $X$ se para qualquer $t\in T$, $$\mathbb{P}(X_t=Y_t)=1$$
\end{dfn}

\vspace{0.5cm}

\noindent Intuitivamente, uma modificação de um processo estocástico $X$ é todo o processo em que o conjunto onde este difere de $X$ tem medida de probabilidade nula. 

\vspace{0.5cm}

\begin{thm}[Kolmogorov-Chentsov]
    Seja  $\{X_t,t \geq0\}$ um processo estocástico tal que existem $K,\alpha,\beta>0$ tal que para todo o $t,s \geq0$
    $$\mathbb{E}[|X_t-X_s|^\alpha]\leq K|t-s|^{1+\beta}.$$

    Então o processo $X$ admite uma modificação contínua $\tilde{X}$.

    Além disso, para qualquer $\gamma \in (0,\beta/\alpha)$ e $T>0$ o processo $\tilde{X}$ é $\gamma$-Holder contínuo em $[0,T]$.
\end{thm}
\begin{proof}
    A prova deste resultado famoso pode ser encontrada em \cite{Kallenberg1997}.
\end{proof}

\vspace{0.5cm}
\noindent A ideia agora é aplicar o teorema anterior ao movimento Browniano fracionário:
\vspace{0.5cm}

\begin{cor}
    O movimento Browniano fracionário admite uma modificação contínua.

    Além disso, para qualquer $\gamma \in (0,H)$ essa modificação é $\gamma$-Holder contínua em qualquer intervalo finito.
\end{cor}

\begin{proof}
    Comecemos por verificar que $\frac{B^H_t-B^H_s}{|t-s|^H}$ é uma variável normal standard. De facto,
    \begin{align*}
        \mathbb{E}[(B^H_t-B^H_s)^2]&=\mathbb{E}[(B_t^H)^2]+\mathbb{E}[(B_s^H)^2]-2\mathbb{E}[B_t^HB_s^H]\\
                                   &=\frac{1}{2}(t^{2H}+t^{2H})+\frac{1}{2}(s^{2H}+s^{2H})-(t^{2H}+s^{2H}-|t-s|^{2H})\\
                                   &=|t-s|^{2H}
    \end{align*}

    Portanto, usando o facto de que para $X \sim \mathcal{N}(0,1)$,
    $$\mathbb{E}[X^k]= \begin{cases}
        \frac{k!}{2^{k/2}(\frac{k}{2})!}, \;\;\;\; \text{n par}\\
        0, \;\;\;\;\;\;\;\;\;\;\;\;\, \text{n ímpar}
                        \end{cases},$$

    temos que $$\mathbb{E}[|B_t^H-B_s^H|^\alpha]=|t-s|^{\alpha H}\mathbb{E}\left[ \left( \frac{|B_t^H-B_s^H|}{|t-s|^H} \right)^\alpha \right]\leq C_\alpha|t-s|^{\alpha H}.$$

    O resultado segue tomando $\beta =\alpha H-1$ e  $ \alpha>1/H$ .

    

\end{proof}



\subsection{Propriedade de Markov, martingale e semimartingale}
Nesta subsecção fixamos um espaço de probabilidade filtrado $(\Omega, \mathcal{F},(\mathcal{F}_t)_{t \geq 0}, \mathbb{P})$.

\begin{dfn}
Um processo estocástico $(X_t)_{t\geq 0}$ diz-se ter a propriedade de Markov se
\begin{equation*}
    P(X_{t+s}\in A|\mathcal{F}_t)=P(X_{t+s} \in A |X_t),
\end{equation*}
para quaisquer $s,t \geq 0$ e $A$ mensurável.
\end{dfn}

\begin{prop}
    Um processo Gaussiano $(X_t)_{t\geq0}$ é de Markov se e só se para quaisquer $0\leq s \leq t \leq u$, 
    $$\operatorname{Cov}(X_s,X_u)=\frac{\operatorname{Cov}(X_s,X_t)\operatorname{Cov}(X_t,X_u)}{\operatorname{Var}(X_t)}$$
\end{prop}
\begin{proof}
    A demonstração pode ser encontrada em \cite{Kallenberg1997}.
\end{proof}

\vspace{1cm}

\begin{thm}
    O movimento Browniano fracionário é de Markov se e só se $H=1/2$.
\end{thm}
\begin{proof}

    Sejam $0<s<t<u$, temos

    \begin{align*}
        \frac{\operatorname{Cov}(B^H_s,B^H_t)\operatorname{Cov}(B^H_t,B^H_u)}{\operatorname{Var}(B^H_t)}&=\frac{1}{4t^{2H}}(t^{2H}+s^{2H}-|t-s|^{2H})(t^{2H}+u^{2H}-|u-t|^{2H})\\
        &=\frac{1}{4t^{2H}} [ t^{4H}+(tu)^{2H}-(t(u-t))^{2H}+(st)^{2H}+(su)^{2H}-s(u-t))^{2H}\\
        &\hspace{0.4cm}+(t(t-s))^{2H}-(u(t-s))^{2H}-((t-s)(u-t))^{2H} ]
    \end{align*}
    

Fazendo $H=\frac{1}{2}$,

\begin{align*}
    \frac{1}{4t} (t^2+tu-tu+t^2+st+su-su+st-t^2+ts-ut+us+tu-t^2-su+st )&=\frac{4st}{4t}\\ 
                                                                        & =s\\
                                                                        &=\min(s,u)\\
                                                                        &=\operatorname{Cov(B^{1/2}_s,B_u^{1/2})}
\end{align*}

Isto é, para $H=1/2$ o \textit{mBf} é de Markov.

Suponha-se agora que $H \neq 1/2$. Para $s=1,t=2,u=3$ temos

\begin{align*}
    \frac{\operatorname{Cov}(B^H_s,B^H_t)\operatorname{Cov}(B^H_t,B^H_u)}{\operatorname{Var}(B^H_t)}&=\frac{1}{4^{H+1}}(16^H+36^H-4^H+4^H+9^H-1-4^H-9^H+1)\\
    &=\frac{16^H+36^H-4^H}{4^{H+1}}=4^{-1}(4^H+9^H-1)\\
\end{align*}

Além disso, $\operatorname{Cov(B_u^H,B_s^H)}=\frac{1}{2}(9^H+1-4^H)$

Finalmente, igualando as duas últimas expressões vem,
$$\frac{1}{2}(9^H+1-4^H)=\frac{4^H+9^H-1}{4} \Leftrightarrow 9^H-3.4^H+3=0,$$

que tem solução única $H=\frac{1}{2}$ em $(0,1)$.
\end{proof}

\vspace{1cm}

\begin{dfn}\label{dfn:martingale}
    Seja $(X_t)_{t\geq0}$ um processo estocástico adaptado à filtração $(\mathcal{F}_t)_{t\geq0}$. Dizemos que $(X_t)_{t \geq 0}$ é uma martingale se
    
    \begin{enumerate}
        \item $\mathbb{E}[|X_t|]<\infty$, para qualquer t
        \item $\mathbb{E}[X_t |\mathcal{F}_s]=X_s$, para quaisquer $s<t$
    \end{enumerate}
\end{dfn}

\vspace{1cm}

\begin{prop}
    Seja $(X_t)_{t \geq0}$ um processo Gaussiano centrado. Então $(X_t)_{t \geq0}$ é uma martingale se e só se para quaisquer $s<t$, $$\operatorname{Cov}(X_t,X_s)=\operatorname{Var}(X_s)$$
\end{prop}

\begin{proof}
    A prova pode ser consultada em \cite{Kallenberg1997}.
\end{proof}

\vspace{1cm}

\begin{thm}
    O movimento Browniano fracionário é uma martingale se e só se $H=1/2.$
\end{thm}
\begin{proof}   
    \begin{align*}
        \operatorname{Cov}(X_t,X_s)=\operatorname{Var}(X_s) \Leftrightarrow \frac{1}{2}(t^{2H}+s^{2H}-|t-s|^{2H})=s^{2H} \Leftrightarrow t^{2H}-(t-s)^{2H}=s^{2H}.       
    \end{align*}

    Se $H=\frac{1}{2}$, então a igualdade anterior reduz-se a $t-(t-s)=s$, a qual é satisfeita para quaisquer $s<t$.\\
    
    Por outro lado se $H \neq\frac{1}{2}$ e se $t=2s$, a expressão anterior é igual a
    $$(2s)^{2H}-(2s-s)^{2H}=s^{2H} \Leftrightarrow 2^{2H}s^{2H}-s^{2H}=s^{2H} \overset{s > 0}{\Leftrightarrow} 2^{2H}-1=1 \Leftrightarrow H=\frac{1}{2}$$
\end{proof}

\vspace{1cm}
De seguida introduzimos alguns conceitos necessários para a definição de semi-martingale.

\vspace{0.5cm}

\begin{dfn}
    Seja $\tau$ uma variável aleatória que toma valores em $T$. Então $\tau$ diz-se um tempo de paragem se 
    $$\{ \tau \leq t\} \in \mathcal{F}_t, \;\;\text{  para qualquer } t \in T$$
\end{dfn}

\vspace{0.2cm}

\begin{dfn}
    Um processo estocástico $\{ (X_t)_{t \geq 0} \}$ diz-se uma martingale local se existe uma sequência $\{\tau_k\}_{k \in \mathbb{N}}$ de tempos de paragem tal que 
    \begin{itemize}
        \item é uma sequência crescente quase certamente, isto é: $\mathbb{P}(\tau_k<\tau_{k+1})=1$;
        \item é uma sequência divergente quase certamente, isto é: $\mathbb{P}(\lim_{k \rightarrow\infty}\tau_k= \infty)=1$;
        \item o processo $X_t^{\tau_k} \coloneq X_{\min\{t,\tau_k\}}$ é uma martingale para todo o $k \in \mathbb{N}$.
    \end{itemize}
\end{dfn}

\vspace{0.2cm}

\begin{dfn}
    Um procesos estocástico $\{ (X_t)_{t \geq 0} \}$ diz-se de Càdlàg se, para qualquer $\omega \in \Omega$
    \begin{itemize}
        \item $\lim_{s \downarrow t} X_s(\omega) = X_t(\omega)$, para qualquer $t \geq0$\\
        \item $X_{t-} \coloneq \lim_{s \uparrow t} X_s(\omega)$ existe para qualquer $t >0$
    \end{itemize}
\end{dfn}

\vspace{0.2cm}

Finalmente chegamos à definição de semi-martingale.

\begin{dfn}
    Um processo $(X_{t})_{t \geq 0}$ diz-se uma semi-martingale se poder ser decomposto como 
    $$X_t=M_t+A_t,$$

    onde $(M_t)_{t \geq 0}$ é uma martingale local e $(A_t)_{t \geq 0}$ é um processo de Càdlàg de variação limitada.

\end{dfn}

\vspace{
0.5cm}

\begin{thm}
 O movimento Browniano fracionário não é uma semi-martingale para $H \neq \frac{1}{2}$
\end{thm}

\vspace{0.5cm}

\begin{lema}
    Se o processo estocástico $(X_t)_{t \geq 0}$ for uma semi-martingale, então $(X_t)_{t \geq 0}$ tem variância quadrática finita.
\end{lema}

\vspace{0.5cm}
    
\begin{proof}{(Teorema 1.8)}

Assuma-se que, para $H \in (0,1)$, $(B^H)_{t \geq 0}$ é uma semi-martingale.

Fixe-se $H\in (0,1)$ e considere-se, para $p>0$ fixo,
\begin{equation*}
    Y_{n,p} \coloneq 2^{n(pH-1)}\sum_{j=1}^{2^n}|B_{j2^{-n}}^H-B_{(j-1)2^{-n}}^H|^p
\end{equation*}

Pela propriedade de autossimilaridade do \textit{mBf} temos 

\begin{equation*}
    Y_{n,p} \overset{\mathcal{D}}{=} 2^{n(pH-1)}\sum_{j=1}^{2^n}2^{{-nHp}}|(B_{j}^H-B_{j-1}^H)|=2^{-n}\sum_{j=1}^{2^n}|(B_{j}^H-B_{j-1}^H)| \eqqcolon \tilde{Y}_{n,p}
\end{equation*}

A sequência $B^H_k-B^H_{k-1}$ é uma sequência i.i.d e portanto ergódica. Logo, pelo Teorema de Ergodic temos que 
\begin{equation*}
    \tilde{Y}_{n,p} \overset{n\rightarrow\infty}{\longrightarrow} \mathbb{E}[|B^H_{1}-B^H_0|^p]=\mathbb{E}[|B_1^H|^p]
\end{equation*}
quase certamente em $L^1$, o que implica convergência em (medida) probabilidade. Como $Y_{n,p} \overset{\mathcal{D}}{=}\tilde{Y}_{n,p}$ temos

\begin{equation*}
    Y_{n,p} \overset{\mathbb{P}}{\longrightarrow}\mathbb{E}[|B_1^H|^p]
\end{equation*}

Finalmente, definimos a p-variância de $B^H$ como 

\begin{equation*}
    V_{n,p}\coloneq\sum_{j=1}^{2^n} |B^H_{j2^{-n}}-B^H_{(j-1)2^{-n}}|^p
\end{equation*}

\vspace{0.5cm}


Observe-se que $V_{n,p}=2^{n(pH-1)}$, logo

\begin{equation*}
    V_{n,p} \overset{\mathbb{P}}{\longrightarrow} 
    \begin{cases}
        0 \;\;\;\;\;\;\;\;, se\;\; pH >1\\
        +\infty\;\;\;\,, se\;\; pH<1
    \end{cases}
\end{equation*}

Ora, para a variância quadrática temos que
\begin{equation*}
    V_{n,2} \overset{\mathbb{P}}{\longrightarrow} 
    \begin{cases}
        0 \;\;\;\;\;\;\;\;, se\;\; H >1/2\\
        +\infty\;\;\;\,, se\;\; H<1/2
    \end{cases}
\end{equation*}

Se $H<1/2$, então $V_{n,2} \overset{\mathbb{P}}{\longrightarrow} +\infty$, o que contradiz o $Lema \;1.1$.

Se $H>1/2$, então $V_{n,2} \overset{\mathbb{P}}{\longrightarrow} 0$. Como, por hipótese $(B^H_t)_{t \geq 0}$ é uma semi-martingale, a sua variação quadrática satisfaz $V_{n,2} \overset{\mathbb{P}}{\longrightarrow}\langle B^H\rangle_1$, logo $\langle B^H \rangle_1=0$, isto é, $(B^H_t)_{t \geq 0}$ é um processo de variação finita. Mas, tomando $p \in (1,1/H)$, tem-se $pH<1$, logo $V_{n,p} \overset{\mathbb{P}}{\longrightarrow} \infty$, pelo que chegamos a uma contradição.

 

\end{proof}

\vspace{1cm}

%\subsection{Representaçoes integrais}

\vspace{5cm}

\subsection{Integração com respeito a \textit{mBf}}
\begin{comment}
Começamos por considerar funções determinísticas. Para tais funções não necessitamos de teoria de integração adicional e podemos definir o integral por aproximação a funções simples da seguinte forma:

Seja $g(x)=\sum_{j=1}^n a_j \mathds{1}_{[t_j,t_{j+1})}$ uma função simples. Definimos 
$$\int g\;dB^H_t \coloneq \sum_{j=1}^na_j[B^H_{j}-B^H_{j+1}]$$

Consideremos um função $f$ não necessariamente simples. Se $\{f_n\}_{n \in \mathbb{N}}$ é uma sucessão de funções simples tal que $f_n \uparrow f$, então definimos 
$$\int f\;dB_t^H \coloneq \lim_{n \rightarrow \infty} \int f_n \; dB^H_t$$


\vspace{0.5cm}
\end{comment}

\hspace{0.4cm} O nosso objetivo nesta secção é dar sentido à seguinte expressão 

$$\int \Delta(t) dB_t^H,$$

\noindent onde a integranda pode ou não ser determinística.

Numa primeira análise, poderia pensar-se em aplicar a teoria de cálculo de Itô ao movimento Browniano fracionário. No entanto, tal tentativa seria problemática, já que o cálculo de Itô está definido apenas para processos que são semi-martingais, uma propriedade que, como acabámos de ver, o movimento Browniano fracionário não satisfaz.
\vspace{0.5cm}

Fixemos $1/2<H<1$. Defina-se $\phi(s,t)=H(2H-1)|t-s|^{2H-2}$. 

\vspace{0.1cm}

\begin{prop}
    A função covariância do movimento Browniano fracionário pode ser escrita como
    $$\frac{1}{2}(s^{2H}+t^{2H}-|t-s|^{2H})=\int_0^t \int_0^s \phi(u,v) dudv$$
\end{prop}
\begin{proof}
     A prova pode ser encontrada em \cite{Inkaya2011}.
\end{proof}


\vspace{0.5cm}

Seja $\mathcal{S}(\mathbb{R})$ o espaço de Schwartz. Defina-se em $\mathcal{S}(\mathbb{R})$ o produto interno definido por

$$<f,g>_{\phi} \coloneq \int_\mathbb{R}\int_\mathbb{R} f(t)g(s)\phi(s,t)dsdt$$

Seja $f: \mathbb{R} \rightarrow \mathbb{R}$ mensurável. Dizemos que $f \in L_\phi^2(\mathbb{R})$ se
$$||f||_\phi^2 =\; <f,f>_\phi\;=\int_\mathbb{R}\int_\mathbb{R} f(s)f(t)\phi(s,t)<\infty$$

Além disso, temos que $(L^2_\phi(\mathbb{R}),<.,.>_\phi)$ é um espaço de Hilbert.

\vspace{0.5cm}


Tome-se $\Omega = \mathcal{S}'(\mathbb{R})$ o dual de $\mathcal{S}(\mathbb{R})$. Para $f \in \mathcal{S}(\mathbb{R})$, a aplicação 
$$f \mapsto \exp \left(-\frac{1}{2}||f||_\phi^2 \right)$$ é positiva definida, então pelo teorema de Bochner-Minlos existe uma medida de probabilidade $\mathbb{P}^H$ nos subconjuntos de Borel de $\Omega$ tal que, para qualquer $f \in \mathcal{S}(\mathbb{R})$ $$\int_\Omega e^{i<\omega,f>} d \mathbb{P}^H(\omega)=e^{-1/2 ||f||_H^2}$$

\vspace{0.5cm}
Vejamos agora como definir o integral com respeito ao movimento Browniano fracionário para funções determinísticas.

Seja $g \in L^2_\phi(\mathbb{R})$ função simples determinística, isto é, $g(x)=\sum_{j=1}^n a_j \mathds{1}_{[t_j,t_{j+1})}$. 
Definimos 
$$\int g\;dB^H_t \coloneq \sum_{j=1}^na_j[B^H_{j+1}-B^H_{j}]$$

Consideremos agora um função $f \in L^2_\phi(\mathbb{R})$ não necessariamente simples. Se $\{f_n\}_{n \in \mathbb{N}}$ é uma sucessão de funções simples em $L^2_\phi(\mathbb{R})$ tal que $f_n \uparrow f$, então definimos 
$$\int f\;dB_t^H \coloneq \lim_{n \rightarrow \infty} \int f_n \; dB^H_t,$$

onde a convergência vale em medida $\mathbb{P}^H$.


\vspace{1cm}

\begin{prop}
    Para $f \in L^2_\phi(\mathbb{R})$ determinística, temos 
    $$\int f \; dB_t^H \sim \mathcal{N}(0,||f||^2_H)$$
\end{prop}
\begin{proof}

    
    Seja $f \in L^2(\mathbb{R})_\phi$ determinística.

    Comecemos por considerar uma função simples, $f_n(x)=\sum_{j=1}^n a_j \mathds{1}_{[t_{j+1},t_j)}(x)$.
    
    Temos,  
    \begin{align*}
        \mathbb{E} \left[ \int f_n\; dB^H_t \right] &= \mathbb{E} \left[\sum_{j=1}^n a_j[B^H_{t_{j+1}}-B^H_{t_j}] \right]\\
                                                    &= \sum_{j=1}^n a_j\mathbb{E}[B^H_{t_{j+1}}-B^H_{t_j}]\\
                                                    &= 0
    \end{align*}

    Portanto, para uma $f$, não necessariamente simples, temos

    \begin{align*}
        \mathbb{E} \left[ \int f\; dB^H_t \right] &= \mathbb{E} \left[  \lim_{n \rightarrow \infty} \int f_n\; dB_t^H \right]\\
                                                  &= \lim_{n \rightarrow \infty} \mathbb{E} \left[ \int f_n \; dB_t^H \right]\\
                                                  &= 0
    \end{align*}


    Além disso, 

    \begin{align*}
        \operatorname{Var} \left[ \int f_n\; dB^H_t \right] &= \mathbb{E}\left[ \left(\int f_n\; dB^H_t \right)^2 \right] \\
                                                            &= \mathbb{E} \left[\left( \sum_{j=1}^na_j[B^H_{j+1}-B^H_{j}] \right)^2 \right]\\
                                                            &=\|f_n\|_\phi^2
    \end{align*}

    Finalmente, para uma $f$ , não necessariamente simples, temos
    \begin{align*}
        \operatorname{Var} \left[ \int f\; dB^H_t \right] &= \operatorname{Var} \left[ \lim_{n \rightarrow \infty} \int f_n dB_t^H \right]\\
                                                          &=\lim_{n \rightarrow \infty} \left[ \int f_n dB_t^H \right]\\
                                                          &=\lim_{n \rightarrow \infty} \| f_n \|_\phi^2\\
                                                          &=\|f\|^2_\phi
    \end{align*}
    
    Falta verificar que $\int f dB_t^H$ é gaussiana. De facto, para cada $n$,
    $$\int f_ndB^H_t=\sum_{j=1}^na_j[B^H_{j+1}-B^H_{j}],$$
    que é uma combinação linear de gaussianas conjuntas. Logo
    $$\int f_ndB^H_t \sim \mathcal{N}(0, \|f_n\|_\phi^2).$$

    Finalmente, como $\int f_ndB^H_t \longrightarrow \int fdB^H_t$, temos que 
    $$\int fdB^H_t \sim \mathcal{N}(0, \|f\|_\phi^2).$$
    
\end{proof}

\vspace{0.5cm}
Seja $L^p(\mathbb{P}^H)=L^p$ o espaço das variáveis aleatórias $X:\Omega \rightarrow \mathbb{R}$ tal que 
$$||X||_{L^p(\mathbb{P}^H)}=\mathbb{E}[|X|^p]^{1/p}<\infty$$
\begin{dfn}
    Define-se o funcional exponencial $\varepsilon: L^2_\phi(\mathbb{R}) \rightarrow L^p$ como
    $$\varepsilon(f) \coloneq exp \left( \int_\mathbb{R} f(t) \;dB^H_t- \frac{1}{2}||f||_H^2\right).$$
\end{dfn}

Observe-se que se $f \in L^2_\phi(\mathbb{R})$, então $\varepsilon(f) \in L^p(\mathbb{P}^H)$, para $p \geq 1$. Portanto, este funcional exponencial permite-nos, a partir de uma função em $L^2_\phi$, obter uma variável aleatória em $L^p$.

\vspace{0.5cm}

\begin{lema}
    A expansão linear de $\{ \varepsilon(f),\;f\in L^2_\phi(\mathbb{R})\}$ é densa em $L^p(\mathbb{P}^H)$ para $p\geq1$.
\end{lema}
\begin{proof}
    A demonstração pode ser consultada em \cite{Duncan2000}.
\end{proof}

\vspace{0.5cm}

%Para $f \in L^2_\phi(\mathbb{R})$ escrevemos $<\omega,f>=\int_\mathbb{R}f(t) \;dB^H(t,\omega)$.

Para podermos enunciar o teorema de expansão, precisamos de uma base ortonormada de $L_\phi(\mathbb{R})$. Essa base é definida com auxílio dos polinómios de Hermite:

\begin{dfn}
    Define-se polinómios de Hermite como
    $$h_n(x)=(-1)^ne^{\frac{x^2}{2}}\frac{d^n}{dx^n}\left(e^{-\frac{x^2}{2}}\right)$$
\end{dfn}
\vspace{0.2cm}
Seja $\mathcal{J}$ o conjunto de todos os multi-índices finitos $\alpha=(\alpha_1,\alpha_2,...)$. Seja ainda $\{e_n\}_{n=1}^\infty$ uma base ortonormada de $L^2_\phi (\mathbb{R})$.
Assim, se $\alpha=(\alpha_1,\alpha_2,...) \in \mathcal{J}$, definimos a base ortonormada 

$$\text{H}_\alpha(\omega) \coloneq h_{\alpha_1}(\xi_i)...h_{\alpha_m}(\xi_m),$$
onde $\xi_i \coloneq <\omega,e_i>$

\vspace{0.2cm}

Podemos então enunciar o teorema de expansão em caos de Wiener: 

\begin{thm}
    Seja $F \in L^2(\mu)$. Então existe constantes $c_\alpha \in \mathbb{R}$ e $\alpha \in \mathcal{J}$, tal que,
    $$F(w)=\sum_{\alpha \in \mathcal{J}}c_\alpha \text{H}_\alpha(w),$$
    onde a convergência vale em $L^2(\mu)$.
    
    \vspace{0.3cm}
    
    Além disso,
    $$||F||^2_\mu=\sum_{\alpha \in \mathcal{J}}\alpha!c_\alpha^2,$$
    onde $\alpha!=\alpha_1!...a_m!$.
\end{thm}

\begin{dfn}
    Define-se 
    $$(S)^*_H \coloneq \left\{G(\omega) = \sum_{\alpha \in \mathcal{J}}c_\alpha H_\alpha(\omega) : ||G||^2_{H,-q}=\sum_{\alpha \in \mathcal{J}}\alpha!c_\alpha^2(2\mathbb{N})^{-\alpha q} < \infty,\;\text{para algum } q \in \mathbb{N}  \right\}$$
\end{dfn}

\begin{dfn}
    Para $F(\omega)=\sum_{\alpha \in \mathcal{J}}a_\alpha \text{H}_\alpha(\omega)$,$G(\omega)=\sum_{\alpha \in \mathcal{J}}b_\alpha \text{H}_\alpha(\omega) \in (S)^*_H$, define-se o produto de Wick de $F$ e $G$ como
    $$(F \lozenge G)(\omega) \coloneq \sum_{\alpha, \beta \in \mathcal{J}}a_\alpha b_\beta \text{H}_{\alpha+\beta}(\omega)$$
\end{dfn}

\begin{prop}
    Sejam $f,g \in L^2_\phi(\mathbb{R})$. Então,
    $$\varepsilon(f) \lozenge \varepsilon(g)=\varepsilon(f+g)$$
\end{prop}

\begin{dfn}
    O \textit{fractional white noise} $W^H_t$ no instante $t$ é definido como 
    $$W_H(t)= \sum_{i=0}^\infty \left[ \int_\mathbb{R}e_i(\nu)\phi(t,\nu)d\nu\right]<\omega, e_i>$$
\end{dfn}

\begin{dfn}
    Se $X : \mathbb{R}  \rightarrow (S)^*_H$ é uma função tal que $X(t) \lozenge W_H(t)$ é integrável em $(S)^*_H$, então definimos
    $$\int_\mathbb{R}X(t)\;dB^H_t \coloneq \int_\mathbb{R} X(t) \lozenge W^H_t dt$$
\end{dfn}
\newpage

\section{Aplicação às finanças - avaliação de opções }

\hspace{0.4cm} Segundo Aristóteles, Tales de Mileto previu uma abundante colheita de azeitonas no próximo verão. Antecipando que tal cenário implicaria uma maior demanda nas máquinas de processamento de azeitonas, Tales procurou tirar partido disso. No entanto, Tales não tinha fortuna suficiente para comprar as tais máquinas. Assim, Tales pagou uma pequena quantia aos donos das máquinas que lhe conferia o direito, mas não a obrigação de as usar no verão. De facto, o verão chegou e com ele uma grande colheita de azeitonas. Tales, que tinha agora o controle sobre as máquinas, alugou-as a um preço superior ao que pagou pelo seu contrato lucrando com a diferença. Este episódio é frequentemente apontado como o primeiro contrato de "opção" documentado da humanidade.

\vspace{0.2cm}

Uma opção é um contrato que concede ao titular o direito, mas não a obrigação, de comprar ou vender uma determinada quantidade de um ativo subjacente por um preço de exercício previamente acordado, numa data de vencimento específica. Denotamos o preço de exercício por $K$ e a data de exercício por $T$, a qual designamos por maturidade.

\vspace{0.2cm}

As opções mais usadas são as opções Europeias e Americanas. As opções Europeias permitem o exercício da opção apenas no instante $T$ (maturidade). Em contraste, uma opção Americana permite o exercício da opção em qualquer instante $t \leq T$.

\vspace{0.2cm}

As opções podem ser de compra, também conhecidas como \textit{call options}, ou de venda, conhecidas como \textit{put options}.
Geralmente, as opções de compra são adquiridas quando se espera que o mercado suba. Denotando por $S_T$ o valor do ativo no instante $T$, temos que a função de retorno de uma opção deste tipo é dada por $(S_T-K)^+$. Por outro lado, as opções de venda são adquiridas quando se espera que o mercado desça. A função de retorno desta opção é dada por $(K-S_T)^+$.

\vspace{0.5cm}

\begin{figure}[H]
    \centering

    \begin{subfigure}{0.48\textwidth}
    \centering
    \begin{tikzpicture}
    \begin{axis}[
        width=\linewidth,
        axis lines=middle,
        xmin=0,xmax=2,
        ymin=-0.2,ymax=1.2,
        xlabel={$S_t$},
        ylabel={Payoff},
        xtick={1},
        xticklabels={$K$},
        ytick=\empty,
    ]

    
    
    % horizontal part
    \addplot[blue, thick] coordinates {(0,0) (1,0)};
    % increasing part
    \addplot[blue, thick] coordinates {(1,0) (2,1)};
    \end{axis}
    \end{tikzpicture}

    \caption{Opção de compra}
    \label{fig:callpayoff}

    \vspace{0.2cm}
    
    $(S_T-K)^+$
\end{subfigure}
    \hfill
\begin{subfigure}{0.48\textwidth}
    \centering
    \begin{tikzpicture}
    \begin{axis}[
        width=\linewidth,
        axis lines=middle,
        xmin=0,xmax=2,
        ymin=-0.2,ymax=1.2,
        xlabel={$S_t$},
        ylabel={Payoff},
        xtick={1},
        xticklabels={$K$},
        ytick=\empty,
    ]
    % decreasing part
    \addplot[red, thick] coordinates {(0,1) (1,0)};
    % horizontal part
    \addplot[red, thick] coordinates {(1,0) (2,0)};
    \end{axis}
    \end{tikzpicture}

    \caption{Opção de venda}
    \label{fig:callpayoff}

    \vspace{0.2cm}
    
    $(K-S_T)^+$
\end{subfigure}

\end{figure}





\begin{comment}
\subsection{Modelo de Black-Scholes-Merton clássico v_0}
Neste modelo fazemos algumas suposições:

\begin{enumerate}
    \item o preço do ativo, $S(t)$, segue um movimento Browniano geométrico, isto é,
\begin{equation}
    dS(t)=\mu S(t)dt+ \sigma S(t) dW(t)
\end{equation}
onde $W(t)$ é um movimento Browniano clássico.

    \item 
    \item 
    \item 

\end{enumerate}



\vspace{0.5cm}

Consideremos um portofólio avaliado, em cada instante $t$ em $X(t)$, constituido por $x_t$ unidades de quantia em dinheiro e $y_t$ unidades de ativo. O portofólio, no instante $t$, tem valor
\begin{equation}
    X(t)=x_te^{rt}+y_tS(t)
\end{equation}
Além disso,
\begin{align}
    dX(t) &= rx_te^{rt}dt + y_tdS(t) \notag\\
    &= rx_te^{rt}dt + y_t(\mu S(t)dt+ \sigma S(t) dW(t)) \notag\\
    &= [rx_te^{rt}+y_t\mu S(t)]dt+y_t\sigma S(t)dW(t)
\end{align}


Por outro lado, pela fórmula de Itô, podemos escrever

\begin{align}
    dX(t,S(t)) &= \frac{\partial X(t,S(t))}{\partial t}dt+\frac{\partial X(t,S(t))}{\partial S(t)}dS(t) +\frac{1}{2}\frac{\partial^2X(t,S(t))}{\partial S(t)^2}dS(t)^2  \notag\\
    &=\frac{\partial X(t,S(t))}{\partial t}dt+\frac{\partial X(t,S(t))}{\partial S(t)}(\mu S(t)dt+ \sigma S(t) dW(t)) +\frac{1}{2}\frac{\partial^2X(t,S(t))}{\partial S(t)^2}(\mu S(t)dt+ \sigma S(t) dW(t))^2  \notag\\
    &= \frac{\partial X(t,S(t))}{\partial t}dt+\frac{\partial X(t,S(t))}{\partial S(t)}(\mu S(t)dt+ \sigma S(t) dW(t)) +\frac{1}{2}\frac{\partial^2X(t,S(t))}{\partial S(t)^2}\sigma^2S^2(t)dt   \notag\\
    &= \left( \frac{\partial X(t,S(t))}{\partial t}+\mu S(t)\frac{\partial X(t,S(t))}{\partial S(t)}+\frac{1}{2}\frac{\partial^2X(t,S(t))}{\partial S(t)^2}\sigma^2S^2(t) \right)dt+ \sigma S(t)\frac{\partial X(t,S(t))}{\partial S(t)}dW(t) \label{eq:important}
\end{align}


\vspace{1cm}

Igualando os coeficientes dos termos deterministicos e aleatórios, vem

\begin{align}
    rx_te^{rt} &=\frac{\partial V}{\partial t}+\frac{1}{2}\frac{\partial^2V}{\partial S^2}\sigma^2 S^2(t)\\
    y_t &=\frac{\partial V}{\partial S}
\end{align}

Substituindo estes valores de $x_t$ e $y_t$ em $(14)$ obtemos

\vspace{0.5cm}

\begin{equation}
    \frac{\partial X}{\partial t} +\frac{1}{2} \sigma^2 S^2 \frac{\partial^2X}{\partial S^2}+rS\frac{\partial X}{ \partial S} -Xr=0
\end{equation}
\end{comment}

\subsection{Modelo de Black-Scholes-Merton clássico}
Neste modelo fazemos algumas suposições:

\begin{enumerate}
    \item o preço do ativo, $S(t)$, segue um movimento Browniano geométrico, isto é,
\begin{equation*}
    dS(t)=\mu S(t)dt+ \sigma S(t) dW(t)
\end{equation*}
onde $W(t)$ é um movimento Browniano clássico.

    \item o mercado é perfeito e completo, isto é, não existem custos de transação; os ativos são infinitamente divisíveis;
    \item não são permitidas oportunidade de arbitrage

\end{enumerate}



\vspace{0.5cm}



Consideremos um portfólio constituído por uma opção e uma quantidade $\Delta$ de ativo, isto é, $$\Pi=C(S,t)-\Delta \times S(t).$$
A ideia é tentar cancelar os fatores de risco do portfólio ao longo do tempo, para isso precisamos de saber como o portfólio evolui ao longo do tempo, isto é, precisamos de conhecer $$d\Pi=dC(S(t),t)-\Delta \times dS(t).$$

Segundo a nossa hipótese de que o ativo segue um movimento Browniano geométrico, sabemos que $dS(t)=\mu S(t)dt+ \sigma S(t)dW(t)$. 

Como o termo $dC(S(t),t)$ depende de um processo estocástico $S(t)$, aplicamos o Lema de Itô:
\begin{equation*}
    dC=\frac{\partial C}{\partial t}dt+\frac{\partial C}{\partial S}dS+\frac{1}{2}\sigma^2 S^2\frac{\partial^2C}{\partial S^2}dt
\end{equation*}

Substituindo e agrupando os termos, temos que 

\begin{align*}
    d\Pi&=\left( \frac{\partial C}{\partial t}dt+\frac{\partial C}{\partial S}dS+\frac{1}{2}\sigma^2 S^2\frac{\partial^2C}{\partial S^2}dt \right)-\Delta \times dS\\
    &=\left( \frac{\partial C}{\partial t} + \frac{1}{2}\sigma^2 S^2\frac{\partial^2C}{\partial S^2}\right)dt+ \left( \frac{\partial C}{\partial S}-\Delta \right)dS
\end{align*}

\vspace{0.5cm}

Portanto, o nosso portfólio é constituído por dois termos: o termo determinístico \textit{dt} e o termo estocástico \textit{dS}. Ora, tomando para $\Delta$, a quantidade de ativos, o valor $\frac{\partial C}{\partial S}$, o termo estocástico anula-se e obtemos um portfólio inteiramente determinístico. Isto é, obtemos um portfólio sem risco (\textit{risk free portfólio}) que rende uma taxa sem risco (\textit{risk free rate}), isto é, 
$$d\Pi=r\Pi dt=r \left( C- \frac{\partial C}{\partial S} \times S \right)dt$$

Igualando as últimas 2 igualdades, obtemos a famosa equação de Black-Scholes

\begin{align*}
    rC- rS\frac{\partial C}{\partial S} =  \frac{\partial C}{\partial t} + \frac{1}{2}\sigma^2 S^2\frac{\partial^2C}{\partial S^2}
\end{align*}

$$\boxed{\frac{\partial C}{\partial t} + \frac{1}{2}\sigma^2 S^2\frac{\partial^2C}{\partial S^2}+rS\frac{\partial C}{\partial S}-rC=0}$$

\vspace{2cm}

\subsection{Risk neutral measure}
Considere-se um contrato financeiro que tem, no instante $T$, retorno aleatório $X$. Qual será o preço "justo" deste contrato num instante $t<T$? Uma primeira resposta natural seria propor que esse preço fosse igual ao valor esperado do tal retorno, descontado ao instante presente. Isto é, $V_t=e^{-r(T-t)}\mathbb{E}[X]$ 

\vspace{0.2cm}

No entanto, esta expressão levanta uma primeira questão: relativamente a que medida de probabilidade deve ser calculado o valor esperado anterior?

\vspace{0.5cm}

Considere-se os seguintes dois investimentos 

\vspace{0.5cm}

\textit{Investimento 1}: Retorno de exatamente 100

\vspace{0.1cm}

\textit{Investimento 2}: Retorno de 
$\begin{cases}
    200 \text{ com probabilidade } \frac{1}{2}\\
    0 \text{ com probabilidade } \frac{1}{2}
\end{cases}$

\vspace{0.25cm}

Note-se que o valor esperado do retorno de ambos os investimentos é o mesmo
\begin{align*}
    \mathbb{E}[I_1]&=100\\
    \mathbb{E}[I_2]&=200 \times \frac{1}{2}+0 =100
\end{align*}

\vspace{0.25cm}

No entanto, é seguro dizer-se que a maioria escolheria o primeiro investimento em detrimento do segundo. Porquê? Porque o segundo envolve maior risco para o investidor. Desta forma, os investidores exigem uma certa recompensa pelo investimento de risco, isto é, requerem que o preço do investimento de risco seja menor que o investimento garantido.\\

\vspace{0.25cm}

Portanto, o preço de um investimento, ao qual está associado um certo risco, será, em geral, diferente do preço atual descontado. 

$$\text{Preço = Retorno esperado (descontado) - Ajuste de Risco}$$

O problema em seguir o diagrama anterior é que cada investimento tem o seu risco associado, e portanto, teríamos que ajustar os preços de cada um individualmente.

\vspace{0.25cm}

A ideia para resolver este problema é inverter os passos. Em vez de descontar o valor esperado e depois ajustar devido ao risco, ajustamos devido ao risco e depois descontamos o valor esperado. Isto é, queremos encontrar uma medida de probabilidade, $\mathbb{Q}$, que considere estes ajustes devido ao risco. Esta tal medida de probabilidade é a chamada \textit{risk neutral measure}.

\vspace{0.25cm}


Matematicamente, queremos construir $\mathbb{Q}$ tal que, para quaisquer $0 \leq t <T$
$$S_t=\mathbb{E^Q}[e^{-(T-t)}S_T|\mathcal{F}_t]$$
o que é equivalente a dizer que, multiplicando ambos os membros por $e^{-rt}$,
$$e^{-rt}S_t=\mathbb{E^Q}[e^{-rT}S_T |\mathcal{F}_t],$$

\vspace{0.25cm}

Isto é, sobre uma medida \textit{risk neutral }, o preço descontado do ativo é uma martingala.

\vspace{0.2cm}

\begin{dfn}
    Uma medida de probabilidade $\mathbb{Q}$ diz-se risk neutral se
    \begin{itemize}
        \item $\mathbb{Q}$ e $\mathbb{P}$ são equivalentes, isto é, $\mathbb{Q}(A)=0$ se e só se $\mathbb{P}(A)=0$
        \item sobre $\mathbb{Q}$ o preço do ativo descontado, $e^{-r(T-t)}S(t)$, é uma martingale
    \end{itemize}
\end{dfn}

\vspace{0.2cm}

\begin{dfn}
    Arbitrage é um portfólio de valor $X(t)$ que satisfaz $X(0)=0$ e para algum $T >0$,
    $$\mathbb{P}\{X(T) \geq 0\}=1, \;\; \mathbb{P}\{ X(T)>0 \}>0$$
\end{dfn}

Intuitivamente, arbitrage é uma oportunidade na qual um investidor começa com capital nulo e em algum tempo posterior $ T$ é certo não ter perdido dinheiro e, portanto, tem uma probabilidade positiva de ter feito dinheiro.


\begin{thm}[Primeiro Teorema Fundamental da Avaliação de Ativos] Se um modelo do mercado tem uma medida de probabilidade risk neutral, então o mercado não admite arbitrage.
\end{thm}
\begin{proof}
A prova pode ser consultada em \cite{Shreve2004}.
\end{proof}

\vspace{0.5cm}

No modelo de Black-Scholes a construção da medida $\mathbb{Q}$ é feita através do Teorema de Girsanov:

\begin{thm}[Girsanov]
    Seja $(W_t)_{t \geq 0} $ um movimento Browniano definido um espaço de probabilidade $(\Omega, \mathcal{F},(\mathcal{F}_t)_{t \geq0}, \mathbb{P})$. Seja $(\Theta)_{t \geq 0}$ um processo adaptado. Defina-se
    \begin{align*}
        &Z(t)=exp \left\{-\int_0^t\Theta(u)dW(u)-\frac{1}{2}\int_0^t\Theta^2(u)du\right\}\\
        &Z=Z(T)\\
        &\tilde{W}(t)=W(t)+\int_0^t\Theta(u)du
    \end{align*}

    Assuma-se que $$\mathbb{E}\left[ \int_0^T \Theta^2(u)Z^2(u)du \right] < \infty.$$ 
    
    Então $\mathbb{E}[Z]=1$ e sobre a medida de probabilidade $\mathbb{Q}$ dada por $$\mathbb{Q}(A)=\int_AZ(\omega)d\mathbb{P}(\omega),$$
    o processo $(\tilde{W}(t))_{t\geq 0}$ é um movimento Browniano. 
\end{thm}
\begin{proof}
    A prova pode ser encontrada em \cite{Shreve2004}.
\end{proof}


\vspace{0.5cm}

Resta agora confirmar que podemos construir, a partir do teorema anterior, uma medida \textit{risk neutral.} Para tal, como visto anteriormente, basta verificar que o valor descontado do preço do ativo é uma martingale sobre a nova medida. Por hipótese,

\begin{align*}
    dS_t=\mu S_tdt+\sigma S_tdW_t
\end{align*}

Escolha-se $\Theta=\frac{\mu-r}{\sigma}$

Segundo o Teorema de Girsanov, $\tilde{W}_t=W_t+\Theta t$ é um movimento Browniano sobre $\mathbb{Q}$. Substituindo,

\begin{align*}
    dS_t&= \mu S_t dt +\sigma S_t d[\tilde{W}-\Theta t]\\
        &= \mu S_tdt+ \sigma S_t [d\tilde{W}-\frac{\mu-r}{\sigma}dt]\\
        &= r S_tdt+ \sigma S_t d\tilde{W}
\end{align*}

Aplicando a fórmula de Itô a $\tilde{S}_t \coloneq e^{-rt}S_t$

\begin{align*}
    d\tilde{S}_t&=e^{-rt}dS_t-re^{-rt}S_tdt\\
                &=e^{-rt}[r S_tdt+ \sigma S_t d\tilde{W}]-re^{-rt}S_tdt\\
                &\sigma S_te^{-rt}d \tilde{W}
\end{align*}

Como o termo de drift é nulo, então $\tilde{S}_t$ é de facto uma martingale.

\vspace{2cm}

\begin{thm}
    A solução da equação de Black-Scholes para uma opção Europeia de compra de preço de exercício $K$ e maturidade $T$ é dada por

    $$c(t,S(t))=S(t)N(d_+)-e^{-r(T-t)}K N(d_-),$$

    onde $$d_{\pm}=\frac{\ln (S(t)/K)+(T-t)(r \pm \frac{1}{2}\sigma^2)}{\sigma  \sqrt{T-t}}$$
\end{thm}

\begin{proof}
    O retorno de uma opção de compra Europeia é dada por $(S(T)-K)^+=max\{S(T)-K,0\}$. Podemos escrever,

    \vspace{0.5cm}

    $$c(t,S(t))=\mathbb{E^Q}[e^{-r(T-t)}(S(T)-K)^+|\mathcal{F}(t)],$$

    onde $\mathbb{Q}$ é a medida \textit{risk neutral } que resulta da aplicação do o teorema de Girsanov com $\Theta=\frac{\mu-r}{\sigma}$. Além disso, resulta que 
    
    $$dS(t)=rS(t)dt+\sigma S(t)d \tilde{W}(t),$$

    Pela fórmula de Itô,

    \begin{align*}
        d[\ln(S(t))]&=\frac{1}{S(t)}dS(t)-\frac{1}{2}\frac{1}{S^2(t)}(dS(t))^2\\
                    &=\frac{1}{S(t)}[rS(t)dt+\sigma S(t)d \tilde{W}(t)]-\frac{1}{2}\frac{1}{S^2(t)}[rS(t)dt+\sigma S(t)d \tilde{W}(t)]^2\\
                    &=\frac{1}{S(t)}[rS(t)dt+\sigma S(t)d \tilde{W}(t)]-\frac{1}{2}\frac{1}{S^2(t)}[\sigma^2S^2(t)dt]\\
                    &=\left(r-\frac{1}{2} \sigma^2 \right)dt+ \sigma d \tilde{W}(t)
    \end{align*}

    Integrando ambos os membros de $t $ a $T$,

    \begin{align*}
        \int_t^T d(\ln S(t)) &= \int_t^T \left( r-\frac{1}{2} \sigma^2 \right) du +\int_t^T \sigma d\tilde{W}(u)\\
                             &=\left(r-\frac{1}{2}\sigma^2\right)(T-t)+\sigma[\tilde{W}(T)-\tilde{W}(t)] =\ln S(T)- \ln S(t)
    \end{align*}

    Finalmente, 

    $$S(T)= S(t)\exp \left\{ \left(r-\frac{1}{2}\sigma^2\right)(T-t)+\sigma[\tilde{W}(T)-\tilde{W}(t)] \right\}$$

    Como $W(T)-W(t) \sim \mathcal{N}(0, T-t)$, podemos escrever

    $$S(T)= S(t)\exp \left\{ \left(r-\frac{1}{2}\sigma^2\right)\tau+\sigma\sqrt{\tau} Z \right\},$$

    onde $Z \sim \mathcal{N}(0,1)$ e $\tau \coloneq T-t$ é o chamado \textit{tempo para a expiração}
Assim,

\begin{align*}
    c(t,x)&=\mathbb{E^Q} \left[ e^{-r \tau} \left( x\exp \left\{ \left(r-\frac{1}{2}\sigma^2\right)\tau+\sigma\sqrt{\tau} Z \right\} -K\right)^+  \right]\\
          &=\frac{1}{\sqrt{2\pi}}\int_{-\infty}^\infty e^{-r \tau} \left( x\exp \left\{ \left(r-\frac{1}{2}\sigma^2\right)\tau+\sigma\sqrt{\tau} z \right\} -K\right)^+ e^{-\frac{1}{2}z^2}dz\\
\end{align*}

A expressão 
$$\left( x\exp \left\{ \left(r-\frac{1}{2}\sigma^2\right)\tau+\sigma\sqrt{\tau} z \right\} -K\right)^+ $$
é positiva se e só se 

\begin{align*}
    \exp \left\{ \left(r-\frac{1}{2}\sigma^2\right)\tau+\sigma\sqrt{\tau} z \right\} >K/x \Leftrightarrow z> \frac{1}{\sigma \sqrt{\tau}} \left( \log(K/x)- \left(r-\frac{1}{2}\sigma^2 \right)\tau \right) :=d_-
\end{align*}


Portanto, 

\begin{align*}
    c(t,x) &= \frac{1}{\sqrt{2\pi}}\int_{d_-}^\infty e^{-r \tau}\left( x\exp \left\{ \left(r-\frac{1}{2}\sigma^2\right)\tau+\sigma\sqrt{\tau} z \right\} -K\right) e^{-\frac{1}{2}z^2} dz\\           
           &=\frac{1}{\sqrt{2\pi}}\int_{d_-}^\infty x\exp \left\{ \sigma \sqrt{\tau}z-\frac{z^2}{2} -\frac{\sigma^2 \tau}{2} \right\} dz - \frac{1}{\sqrt{2\pi}} \int_{d_-}^\infty e^{-r\tau}Ke^{-\frac{1}{2}z^2}dz\\
           &=\frac{x}{\sqrt{2\pi}}\int_{-\infty}^{d_-} \exp \left\{ -\sigma \sqrt{\tau} y- \frac{y^2}{2} -\frac{\sigma^2\tau}{2}  \right\}dy +e^{-r\tau} K \int_{-\infty }^{d_-} \frac{1}{\sqrt{2\pi}}e^{-\frac{1}{2}z^2}dz\\
           &=x\int_{\infty}^{d_-}\frac{1}{\sqrt{2\pi}}\exp\left\{ -\frac{1}{2} \left ( y+\sigma\sqrt{\tau} \right)^2 \right\}dy + e^{-r\tau} K \int_{-\infty }^{d_-} \frac{1}{\sqrt{2\pi}}e^{-\frac{1}{2}z^2}dz\\
           &=x\int_{\infty}^{d_-+\sigma \sqrt{\tau}}\frac{1}{\sqrt{2\pi}}e^{-\frac{1}{2}w^2}dw + e^{-r\tau} K \int_{-\infty }^{d_-} \frac{1}{\sqrt{2\pi}}e^{-\frac{1}{2}z^2}dz\\
           &=xN(d_+)+e^{-r \tau}KN(d_-)            
\end{align*}
\end{proof}

\vspace{2cm}

\subsection{Modelo de Black-Scholes-Merton fracionário}
De um modo totalmente análogo ao movimento Browniano, podemos derivar a equação de Black-Scholes fracionário, isto é, quando se assume que o ativo segue um movimento Browniano fracionário geométrico. 
\begin{thm}   

O preço de uma contrato financeiro sobre um ativo de payoff $f(S(T))$ é dado por $V(t,S(t))$, onde $V(t,S)$ é a solução da EDP:

\begin{align*}
    &\frac{\partial V}{\partial t}+H \sigma^2 t^{2H-1}S^2 \frac{\partial^2V}{\partial S^2}+rS\frac{\partial V}{\partial S}-rV=0\\
    &V(T,S)=f(S)
\end{align*}
\end{thm}

\begin{proof}
    Analogamente ao modelo Black-Scholes clássico, consideramos um portfólio constituído por uma opção e uma quantidade $\Delta$ de ativos,isto é, o portfólio, no instante $t$, tem valor $$\Pi(t)=V(t,S(t))-\Delta S(t)$$

    Usando a fórmula fracionária de Ito, vem que 

    \begin{align*}
        d\Pi(t)&=dV(t,S(t))-\Delta dS(t)\\
               &=\left( \frac{\partial V}{\partial t} +H \sigma^2 t^{2H-1}S^2 \frac{\partial^2V}{\partial S^2} -\Delta \delta S \right)dt+ \left( \sigma S\frac{\partial V}{\partial S} -\Delta \sigma S \right) dB^H(t)
    \end{align*}

    Mais uma vez, pretendemos um portfólio  sem risco. Assim, tomamos $\Delta=\frac{\partial V}{\partial S}$. Logo 
    $$d \Pi(t)=r \Pi(t)dt$$

    Finalmente, igualando as últimas duas expressões vem 

    $$\frac{\partial V}{\partial t}+H \sigma^2 t^{2H-1}S^2 \frac{\partial^2V}{\partial S^2}+rS\frac{\partial V}{\partial S}-rV=0$$
    
\end{proof}



\vspace{1cm}


\begin{thm}\label{thm:call}
    O preço de uma opção de compra Europeia de \textit{strike price} $K$ e maturidade $T$, no instante $t$ é dada por:

    $$C(t,S(t))=S(t)N(d_+)-Ke^{-r(T-t)}N(d_-),$$

    onde $$d_+=\frac{\ln(S(t)/K)+r(T-t)+\frac{1}{2}\sigma^2(T^{2H}-t^{2H})}{\sigma \sqrt{T^{2H}-t^{2H}}}$$
    \hspace{0.5cm} e    $$d_-=\frac{\ln(S(t)/K)+r(T-t)-\frac{1}{2}\sigma^2(T^{2H}-t^{2H})}{\sigma \sqrt{T^{2H}-t^{2H}}}$$
\end{thm}
\begin{proof}
    A prova é análoga ao caso do modelo  clássico e pode ser consultada em \cite{Inkaya2011}.
\end{proof}

\vspace{0.5cm}

Observe-se que, avaliando a opção em $t=0$ com maturidade $T=1$, temos que $T^{2H}-t^{2H}=1$

\vspace{0.5cm}

Note-se que o comprador da opção anterior, tem ganhos potênciais ilimitados, enquanto que as perdas estão limitadas. De facto, se o preço do ativo na maturidade $T$ for menor que $K$, $S_T<K$, então a opção não é exercida e as perdas são nulas. Por outro lado, se $S_T>K$, então os ganhos são dados por $S_T-K$ que pode ser arbitrariamente grande, dependendo do valor de $S_T$. Deste modo, faz sentido que um aumento da volatilidade conduza, em geral, a um aumento do preço da opção.\\

No modelo de Black-Scholes fracionário, este comportamento é refletido pelo termo a volatilidade $\sigma \sqrt{T^{2H}-t^{2H}}$.

\vspace{2cm}

\section{Simulações numéricas}
Nesta secção apresentamos algumas simulações de resultados apresentados até agora.

\vspace{0.5cm}

\subsection{Preço de opção de compra Europeia}
Verifiquemos como o preço da opção de compra dada pelo Teorema~\ref{thm:call} varia com o índice de Hurst e com a maturidade $T$. Por simplicidade tomou-se $t=0$, pelo que  $T-t=T$.

\subsubsection{Dependência relativamente ao parâmetro de Hurst}
\begin{figure}[H]
    \centering
    \includegraphics[width=1\textwidth]{HvsC.pdf}
\end{figure}


Uma primeira observação consiste na curva para $T=1$ ser uma reta horizontal. Isto é, o preço da opção, para $T=1$, não depende de $H$. Isto é explicado pelo facto de para $T=1$ e $t=0$, o termo $T^{2H}-t^{2H}=1$ e por este ser o único termo na expressão do preço da opção que dependia de $H$.
\vspace{0.2cm}

\noindent Para maturidades menores que $1$, o preço da opção diminui ao longo de $H$. Isto porque, para $T<1$, o termo $T^{2H}$ é decrescente enquanto $H$ cresce.

\vspace{0.2cm}

\noindent Por outro lado, para maturidades superiores a $1$, o preço da opção aumenta ao longo de $H$. Mais uma vez, para $T>1$, o termo $T^{2H}$ é crescente ao longo de $H$. Além disso, para maturidades $T>>1$, em particular para $T=5$, o termo $T^{2H}$ torna-se muiot mais sensível a alterções de $H$.

\vspace{0.5cm}

\subsubsection{Dependência relativamente à maturidade}

\begin{figure}[H]
    \centering
    \includegraphics[width=1\linewidth]{TvsC.pdf}
\end{figure}

Em primeiro lugar, note-se que todas as curvas se intersetam em $T=1$. Mais uma vez, para $T=1$ (e $t=0$), temos que $T^{2H}-t^{2H}=1$, e portanto, o preço da opção torna-se independente de $H$. 

\vspace{0.2cm}

Para $T<1$ observa-se que valores mais elevados de $H$ conduzem a preços inferiores da opção. Por outro lado, para $T>1$, observa-se o oposto: o preço da ação aumenta com o parâmetro de Hurst. Este comportamento está, mais uma vez, relacionado com o facto de a volatilidade efetiva do modelo ser dada por $$\sigma \sqrt{T^{2H}-t^{2H}}$$ a qual diminui com $H$ para $T<1$ e aumenta com $H$ para $T>1$.

\vspace{0.2cm}

Por fim, observa-se que as curvas divergem progressivamente à medida que a maturidade aumenta. Isto evidência que o impacto do parâmetro de Hurst é reduzido para contractos de curta maturidade. Um comportamento que está relacionado com a limitação das perdas e com a ilimitação dos ganhos do titular da opção, como já foi explicado anteriormente.


\subsubsection{Compração com o modelo de Black-Scholes para o movimento Browniano}

\begin{figure}[H]
    \centering
    \includegraphics[width=1\linewidth]{3222b11e-c848-4af2-8129-d72270ab8214.pdf}
\end{figure}

Mais uma vez, começa-se por observar que para $T=1$ todas as curvas se intersetam em 0. Isto porque, como já vimos, para $T=1$ o preço da opção é independente do parâmetro de Hurst e portanto, o preço das opções em ambos os modelos (clássico e fracionário) será o mesmo.

Para maturiadades inferiores a um ano, verifica-se que a diferença $C_H-C_{BS}$ é negativa quando $H>0.5$ e positiva quando $H<0.5$. Isto significa que, para contratos de curta maturidade, o modelo clássico de Black-Scholes sobreestima o preço da opção quando existe persistência e subestima-o quando existe antipersistência. Para maturidades superiores a um ano observa-se extamente o comportamento oposto.

Finalmente, verifica-se que o módulo da diferença aumenta progressivamente com a maturidade. Este comportamento evidencia que as discrepâncias entre os dois modelos são praticamente negligenciáveis para contratos de curta duração, tornando-se significativamente mais relevantes à medida que a maturidade aumenta.



\newpage


\begin{thebibliography}{99}

\bibitem{Inkaya2011}
A. Inkaya,
\textit{Option Pricing with Fractional Brownian Motion},
M.Sc. Thesis, Department of Financial Mathematics,
Middle East Technical University, Ankara, Turkey, 2011.

\bibitem{Necula2002}
C. Necula,
``Option Pricing in a Fractional Brownian Motion Environment,''
Working Paper,
Academy of Economic Studies, Bucharest, Romania, 2002.

\bibitem{Nourdin2012}
I. Nourdin,
\textit{Selected Aspects of Fractional Brownian Motion},
Bocconi \& Springer Series, Vol. 4,
Springer, 2012.

\bibitem{Sznitman2013}
A.-S. Sznitman,
\textit{Brownian Motion and Stochastic Calculus},
Lecture Notes, Spring Term 2013.

\bibitem{Kallenberg1997}
O. Kallenberg,
\textit{Foundations of Modern Probability},
Springer, New York, 1997.

\bibitem{Rostek2009}
S. Rostek,
\textit{Option Pricing in Fractional Brownian Markets},
Lecture Notes in Economics and Mathematical Systems, Vol. 622,
Springer, 2009.

\bibitem{Shevchenko2014}
G. Shevchenko,
``Fractional Brownian Motion in a Nutshell,''
arXiv:1406.1956, 2014.

\bibitem{Elliott2003}
R. J. Elliott and J. van der Hoek,
``A General Fractional White Noise Theory and Applications to Finance,''
\textit{Mathematical Finance},
Vol. 13, No. 2, pp. 301--330, 2003.

\bibitem{VilelaMendes2012}
R. Vilela Mendes, M. J. Oliveira and A. M. Rodrigues,
``The Fractional Volatility Model: No-Arbitrage, Leverage and Completeness,''
2012.

\bibitem{Shreve2004}
S. E. Shreve,
\textit{Stochastic Calculus for Finance II: Continuous-Time Models},
Springer Finance Textbooks,
Springer, New York, 2004.

\bibitem{RevuzYor1999}
D. Revuz and M. Yor,
\textit{Continuous Martingales and Brownian Motion},
3rd ed.,
Grundlehren der mathematischen Wissenschaften, Vol. 293,
Springer, Berlin, 1999.

\bibitem{Duncan2000}
T. E. Duncan, Y. Hu and B. Pasik-Duncan,
``Stochastic Calculus for Fractional Brownian Motion I. Theory,''
\textit{SIAM Journal on Control and Optimization},
Vol. 38, No. 2, pp. 582--612, 2000.

\end{thebibliography}

\end{document}
